% !TeX root = ../../book.tex
\subsection*{递归定义}

在 \Crefrange{cqRecursiveDefinitionsOfArithmeticOperationsBegin}{cqRecursiveDefinitionsOfArithmeticOperationsEnd} 中，直接使用加法、乘法和指数的递归定义来证明给出的方程。

\begin{chapex}
\label{cqRecursiveDefinitionsOfArithmeticOperationsBegin}
$1+3=4$
\end{chapex}

\begin{chapex}
$0+5=5$
\end{chapex}

\begin{chapex}
$2 \cdot 3 = 6$
\end{chapex}

\begin{chapex}
$0 \cdot 5 = 0$
\end{chapex}

\begin{chapex}
\label{cqRecursiveDefinitionsOfArithmeticOperationsEnd}
$2^3=8$
\end{chapex}

\begin{chapex}
给出 $n \in \mathbb{N}$ 上新量词 $\exists^{=n}$ 的递归定义，其中给定集合 $X$ 和谓词 $p(x)$，逻辑公式 $\exists^{=n} x \in X,~ p(x)$ 表示`恰好存在 $n$ 个元素 $x \in X$ 使得 $p(x)$ 为真'。即定义 $\exists^{=0}$，然后根据 $\exists^{=n}$ 定义 $\exists^{=n+1}$。
\hintlabel{cqRecursiveDefinitionOfExistsExactlyNQuantifier}{%
例如，如果 $X$ 中恰好有 $3$ 个元素使得 $p(x)$ 为真，则意味着存在某个 $a \in X$ 使得 $p(a)$ 为真，且除 $a$ 以外还有恰好两个元素 $x \in X$ 使得 $p(x)$ 为真。
}
\end{chapex}

\begin{chapex}
使用二项式系数的递归定义 (\Cref{defBinomialCoefficientRecursive}) 直接证明 $\dbinom{4}{2} = 6$。
\end{chapex}

\begin{chapex}
\begin{enumerate}[(a)]
\item 求 $41!$ 的十进制展开中尾数 $0$ 的数量。
\item 求 $41!$ 的二进制展开中尾数 $0$ 的数量。
\end{enumerate}
\hintlabel{cqTrailingZerosInFactorial}{%
自然数 $n$ 的 base-$b$ 展开中尾数零的数量是使 $b^r$ 整除 $n$ 的最大自然数 $r$。$10$ 出现在 $41!$ 有多少次？ $2$ 出现在 $41!$ 有多少次？
}
\end{chapex}

\begin{chapex}
设 $N$ 为集合，设 $z \in N$ 并设 $s : N \to N$。证明 $(N, z, s)$ 为自然数（在 \Cref{defNotionOfNaturalNumbers} 的意义上）当且仅当，对于每一个集合 $X$，每一个元素 $a \in X$ 和每一个函数 $f : X \to X$，都有一个特定函数 $h : N \to X$ 满足 $h(z)=a$ 且 $h \circ f = s \circ h$。
\end{chapex}

\subsection*{用归纳法证明}

\begin{chapex}
\label{cqNPowerFiveLessThanFivePowerN}
找到满足 $n^5 < 5^n$ 的所有自然数 $n$。
\end{chapex}

\begin{chapex}
证明对于所有实数 $x \ge -1$, $(1+x)^{123\;456\;789} \ge 1 + 123\;456\;789\;x$。
\hintlabel{cqBernoullisInequality}{%
通过在 $x$ 上应用归纳法来证明这一点仅能证明对于\textit{整数} $x \ge -1$ 成立，而非\textit{实数} $x \ge -1$。尝试通过对不同变量进行归纳来证明更普遍的事实，然后将其作为特殊情况进行推导。（你\textit{真的}认为自然数 $123\;456\;789$ 有什么特别之处吗？）
}
\end{chapex}

\begin{chapex}
设 $a \in \mathbb{N}$ 并假设 $a$ 的十进制展开的最后一位数字是 $6$。证明对于所有 $n \ge 1$, $a^n$ 的十进制展开的最优一位数字为 $6$。
\end{chapex}

\begin{chapex}
设 $a,b \in \mathbb{R}$，并设 $a_0, a_1, a_2, \dots$ 为对于所有 $n \ge 1$ 满足 $a_0 = a$, $a_1=b$ 且 $a_n = \dfrac{a_{n-1} + a_{n+1}}{2}$ 的序列。证明对于所有 $n \in \mathbb{N}$, $a_n = a + (b-a) n$。
\end{chapex}

\begin{chapex}
设 $f : \mathbb{R} \to \mathbb{R}$ 为函数，对于所有 $x,y \in \mathbb{R}$ 满足 $f(x+y) = f(x) + f(y)$。
\begin{enumerate}[(a)]
\item 用归纳法证明对于所有 $n \in \mathbb{N}$ 存在实数 $a$ 使得 $f(n) = an$。
\item 对于所有 $n \in \mathbb{Z}$ 推导出 $f(n) = an$。
\item 对于所有 $x \in \mathbb{Q}$ 推导出 $f(x) = ax$。
\end{enumerate}
\hintlabel{cqHomomorphismFromRPlusToRPlus}{%
要找到 (a) 中 $a$ 的值，尝试将 $x$ 和 $y$ 的某些小数值代入方程 $f(x+y)=f(x)+f(y)$。
对于 (b)，使用对于所有 $n \in \mathbb{N}$, $n + (-n) = 0$。
对于 (c)，考虑 $b f(\frac{a}{b})$ 其中 $a,b \in \mathbb{Z}$ 且 $b \ne 0$。
}
\end{chapex}

\begin{chapex}
设 $f : \mathbb{R} \to \mathbb{R}$ 为函数，对于所有 $x,y \in \mathbb{R}$ 满足 $f(0) > 0$ 且 $f(x+y)=f(x)f(y)$。证明存在某个正实数 $a$ 对于所有\textit{有理}数 $x$ 满足 $f(x)=a^x$。
\hintlabel{cqHomomorphismFromRPlusToRTimes}{%
跟 \Cref{cqHomomorphismFromRPlusToRPlus} 类似，首先证明 $x \in \mathbb{N}$（通过归纳法），然后证明 $x \in \mathbb{Z}$，最后证明 $x \in \mathbb{Q}$。
}
\end{chapex}

\begin{chapex}
设 $a,b \in \mathbb{Z}$。证明对于所有 $n \in \mathbb{N}$, $b-a$ 整除 $b^n-a^n$。
\end{chapex}

\begin{chapex}
用归纳法证明对于所有 $n \in \mathbb{N}$, $7^n - 2 \cdot 4^n + 1$ 能被 $18$ 整除。
\hintlabel{cqInductionEmbeddedInIductionStep}{%
对于所有 $n \in \mathbb{N}$, $7^{n+1} - 2 \cdot 4^{n+1} + 1 = 4(7^n - 2 \cdot 4^n + 1) + 3(7^n - 1)$。你可能需要在归纳步骤中分别用归纳法证明。
}
\end{chapex}

\subsection*{微积分中的一些例子}

\Crefrange{cqCalculusBegin}{cqCalculusEnd} 假设你熟悉微分和积分的基本技巧。

\begin{chapex}
\label{cqCalculusBegin}
证明对于所有 $n \in \mathbb{N}$, $\dfrac{d^n}{dx^n}(xe^x) = (n+x)e^x$。
\end{chapex}

\begin{chapex}
找到并证明对于所有 $n \in \mathbb{N}$, $\dfrac{d^n}{dx^n}(x^2e^x)$ 的表达式。
\end{chapex}

\begin{chapex}
设 $f$ 和 $g$ 为可微函数。证明对于所有 $n \in \mathbb{N}$,
\[ \dfrac{d^n}{dx^n}(f(x)g(x)) = \sum_{k=0}^n \dbinom{n}{k} f^{(k)}(x) g^{(n-k)}(x) \]
其中符号 $h^{(r)}$ 表示 $h$ 的 $r^{\text{th}}$ 导数。
\end{chapex}

\begin{chapex}
找到并证明对于所有 $n \in \mathbb{N}$, $\log(x)$ 的 $n^{\text{th}}$ 导数的表达式。
\end{chapex}

\begin{chapex}
证明对于所有 $\theta \in \mathbb{R}$ 和 $n \in \mathbb{N}$, $(\cos \theta + i \sin \theta)^n = \cos (n\theta) + i \sin (n\theta)$。
\hintlabel{cqDeMoivre}{%
你需要使用角度加法公式
\[ \sin(\alpha \pm \beta) = \sin\alpha\cos\beta \pm \cos\alpha\sin\beta ~~\text{和}~~ \cos(\alpha \pm \beta) = \cos \alpha \cos \beta \mp \sin \alpha \sin \beta \]
}
\end{chapex}

\begin{chapex}
\label{cqCalculusEnd}
\begin{enumerate}[(a)]
\item 证明对于所有 $n \in \mathbb{N}$, $\displaystyle \int_0^{\frac{\pi}{2}} \sin^{n+2}(x)~dx = \dfrac{n+1}{n+2} \int_0^{\frac{\pi}{2}} \sin^n(x)~dx$。
\item 用归纳法证明对于所有 $n \in \mathbb{N}$, $\displaystyle \int_0^{\frac{\pi}{2}} \sin^{2n}(x)~dx = \dfrac{\pi}{2^{2n+1}} \dbinom{2n}{n}$。
\item 找到并证明对于所有 $n \in \mathbb{N}$, $\displaystyle \int_0^{\frac{\pi}{2}} \sin^{2n+1}(x)~dx$ 的表达式。
\end{enumerate}
\hintlabel{cqIntegralOfSinToTheN}{%
对于 (a) 分块使用积分；不需要归纳法。
对于 (b) 和 (c)，使用 (a) 的结论。
}
\end{chapex}

\subsection*{判断题}

\tfquestiontext{cqInductionTFBegin}{cqInductionTFEnd}

\begin{chapex} % True
\label{cqInductionTFBegin}
每个阶乘都是正数。
\end{chapex}

\begin{chapex} % False
每个二项式系数都是正数。
\end{chapex}

\begin{chapex} % False
每个强归纳法证明都可以通过仅改变归纳假设变成弱归纳法证明。
\end{chapex}

\begin{chapex} % True
每个弱归纳法证明都可以通过仅改变归纳假设变成强归纳法证明。
\end{chapex}

\begin{chapex} % True
给定集合 $X$，满射 $f : \mathbb{N} \to X$ 和逻辑公式 $p(x)$，其中自由变量 $x \in X$，如果 $p(f(0))$ 为真且 $\forall n \in \mathbb{N},\, (p(f(n)) \Rightarrow p(f(n+1)))$ 为真，则 $\forall x \in X,\, p(x)$ 为真。
\end{chapex}

\begin{chapex} % False
给定某个逻辑表达式 $p(x)$，其中自由变量 $x \in \mathbb{N}$，如果对于所有 $x \in \mathbb{N}$ 存在某个 $y \le x$ 使得 $p(x)$ 为假，则对于所有 $x \in \mathbb{N}$, $p(x)$ 为假。
\end{chapex}

\begin{chapex} % False
非负有理数集合 $\mathbb{Q}^{\ge 0}$ 的每个非空子集都存在最小元素。 
\end{chapex}

\begin{chapex} % False
整数集 $\mathbb{Z}$ 的每个非空子集都存在最小元素。
\end{chapex}

\begin{chapex} % True
\label{cqInductionTFEnd}
负整数集合 $\mathbb{Z}^{-}$ 的每个非空子集都存在最大元素。
\end{chapex}

\subsection*{可能性问题}

\asnquestiontext{cqInductionASNBegin}{cqInductionASNEnd}

\begin{chapex} % Sometimes
\label{cqInductionASNBegin}
设 $n,k,\ell \in \mathbb{Z}$。则 $\dbinom{n}{k+\ell} = \dbinom{n}{k} + \dbinom{n}{\ell}$。
\end{chapex}

\begin{chapex} % Never
设 $m,n \ge 2$。则 $(m+n)! = m!+n!$。
\end{chapex}

\begin{chapex} % Always
\label{cqInductionASNEnd}
设 $p(x)$ 为逻辑表达式，其中自由变量 $x \in \mathbb{Z}$，假设 $p(0)$ 为真，且对于所有 $n \in \mathbb{Z}$，如果 $p(n)$ 为真，则 $p(n+1)$ 和 $p(n-2)$ 为真。则 $\forall n \in \mathbb{Z},\, p(n)$ 为真。
\end{chapex}